\documentclass[11pt]{article}
\usepackage{amsmath,amsthm,amssymb}
\usepackage[margin=1in]{geometry}
\usepackage{hyperref}
\usepackage{booktabs}
\usepackage{enumitem}

\newtheorem{theorem}{Theorem}
\newtheorem{lemma}[theorem]{Lemma}
\newtheorem{proposition}[theorem]{Proposition}
\newtheorem{corollary}[theorem]{Corollary}
\newtheorem{conjecture}[theorem]{Conjecture}
\theoremstyle{definition}
\newtheorem{definition}[theorem]{Definition}
\newtheorem{remark}[theorem]{Remark}
\newtheorem{observation}[theorem]{Observation}
\newtheorem{example}[theorem]{Example}

\newcommand{\Hq}{H_q}
\newcommand{\Sn}{S_n}
\newcommand{\Vlam}{V_\lambda}
\newcommand{\Bdec}{B_{\ell+1}^{(\lambda)}}
\newcommand{\Om}{\Omega}
\newcommand{\hatl}{\widehat\lambda}
\newcommand{\supp}{\mathrm{supp}}
\newcommand{\SYT}{\mathrm{SYT}}
\newcommand{\Rp}{{R'}}

\title{Extended sign-kill for $\hatl=(5,2)$: \\
       four-piece $E^+_{n-1}$ decomposition, two hook branches}
\author{Clio}
\date{14 May 2026 (very late prove session, tenth pass)}

\begin{document}
\maketitle

\begin{abstract}
Let $\lambda = (5, 2, 1^q) \vdash n = q + 7$ with $q \ge 5$,
$\ell = q + 2 = n - 5$, and
$B = \Bdec\Vlam = \Rp_{\ell+1}\Rp_{\ell+2}\Rp_{\ell+3}\Rp_{\ell+4}\Rp_{\ell+5}\,\Vlam$
(five $R'$-blocks).  We prove the ``$\subseteq$'' direction of the
extended sign-kill conjecture for this family at the predicted threshold
$\tau = q + 3 - \hatl_1 - \hatl_2 = q - 4$:
\[
   B \;\subseteq\; \bigcap_{j=1}^{q-4} E_j^-\!\mid_{\Vlam}.
\]
At $q \in \{0, 1, 2, 3, 4\}$ the predicted threshold is non-positive,
so the containment is vacuous.

The shape $\hatl = (5, 2)$ has $r = 2$ length-preserving corners, so the
inner $E^+_{n-1}$ decomposition has four pieces (three corner-pair
$\Phi$-pieces and one same-row $\Phi_h$-piece) --- the same total count as
the eighth-pass $\hatl = (4, 3)$ proof.  The new combinatorial feature:
\emph{two} of the four pieces are hooks (vs.\ one hook in $(4, 3)$).
Specifically, the contributing branches are $\mu_A = (4, 2, 1^{q-1})$
(non-hook, $\hatl = (4, 2)$) and $\mu_B = (5, 1^q)$ (hook), and the
vanishing branches are $\mu_C = (4, 1^{q+1})$ (hook) and $\mu_h = (3, 2,
1^q)$ (non-hook, $\hatl = (3, 2)$).  The prefix bound on the $\mu_B$
piece comes from the May-14 fifth-pass hook support lemma at $k = 5$,
matching $q - 3$ to the recursive $(4, 2)$ prefix bound on the $\mu_A$
piece.  Dimensions match: $\dim B = f^{(4,1)} = 4 = f^{(3, 1)} + f^{(4)}
= \dim B^{(\mu_A)}_{\ell+1} + \dim B^{(\mu_B)}_{\ell+1}$.

This completes the $\hatl_1 + \hatl_2 = 7$ row of the recursive ladder.
\end{abstract}

\section{Setup}\label{sec:setup}

We use the conventions of the May-13--May-14 series.  $\Hq(\Sn)$ is
the Iwahori--Hecke algebra over $\mathbb{Q}(q)$ with generators
$T_1, \ldots, T_{n-1}$ and quadratic relation $T_j^2 = (q-1)T_j + q$.
$\Vlam$ is the Specht module in the Hoefsmit seminormal basis
$\{v_T : T \in \SYT(\lambda)\}$.

For $1 \le j \le n-1$, set $S_j := T_j + 1$.  Write
\[
   E_j^+ \;:=\; \ker(T_j - q)\!\mid_{\Vlam}
   \;=\; \mathrm{Im}(S_j\!\mid_{\Vlam}),
   \qquad
   E_j^- \;:=\; \ker(S_j)\!\mid_{\Vlam}.
\]
For $k \ge 1$, let $\Rp_k := S_{k-1}\,S_{k-2}\cdots S_1$ with the
convention $\Rp_1 = 1$.  For $\lambda \vdash n$ with $\ell = \ell(\lambda)$,
set
\[
   \Om \;=\; \Om^{(\lambda)} \;:=\; \Rp_{\ell+1}\Rp_{\ell+2}\cdots \Rp_n,
   \qquad
   B \;:=\; \Om \cdot \Vlam \;=\; \Bdec \Vlam.
\]

For $\lambda = (5, 2, 1^q)$, $n = q + 7$ and $\ell = q + 2$, so
$\Om = \Rp_{q+3}\Rp_{q+4}\Rp_{q+5}\Rp_{q+6}\Rp_{q+7} =
\Rp_{n-4}\Rp_{n-3}\Rp_{n-2}\Rp_{n-1}\Rp_n$
has \emph{five} $\Rp$-blocks.  By the May-13 multi-corner-dimension
paper, $\dim B = f^{\lambda^{(\ge 2)}} = f^{(4, 1)} = 4$ in the stable
regime $q \ge 2$.

\paragraph{Removable corners of $\lambda$.}
The removable corners are
\[
   c_1 = (1, 5), \qquad c_2 = (2, 2), \qquad c_3 = (q+2, 1).
\]
$\lambda$ has two length-preserving corners ($c_1, c_2$) and one
column-$1$ bottom corner ($c_3$), an $r = 2$ shape.

\paragraph{Identity prefix and seminormal support.}  For an SYT $T$
of any partition $\mu$, the \emph{identity prefix length} is
\[
   p(T) \;:=\; \max\{\,P \ge 0 : T(i, 1) = i \text{ for all }
   1 \le i \le P\,\}.
\]
For $W \subseteq V_\mu$,
$\supp(W) := \{\,T \in \SYT(\mu) : \exists w \in W,\ [v_T]\,w \ne 0\,\}$.

\section{Input from prior papers}\label{sec:input}

\begin{proposition}[Support rule, May-14 fourth pass]\label{prop:support-rule}
Let $W \subseteq V_\nu$ be any subspace, $1 \le j \le |\nu|-1$.
If every $T \in \supp(W)$ has $j$ and $j+1$ in the same column of $T$,
then $W \subseteq E_j^-\!\mid_{V_\nu}$.
\end{proposition}

\begin{remark}\label{rem:prefix-implies-Ej-}
If every $T \in \supp(W)$ has $p(T) \ge P$, then for every
$1 \le j \le P - 1$ the entries $j, j+1$ sit at column-$1$ rows
$j, j+1$, sharing column~$1$.  By Proposition~\ref{prop:support-rule},
$W \subseteq E_j^-$ for $1 \le j \le P - 1$.
\end{remark}

\begin{theorem}[Extended sign-kill at $\hatl = (3, 2)$, May-14 fifth pass]\label{thm:32-prefix}
For $\lambda' = (3, 2, 1^{q'}) \vdash n' = q' + 5$ with $q' \ge 3$,
every $T' \in \supp\bigl(B^{(\lambda')}_{q'+3} V_{\lambda'}\bigr)$
satisfies $p(T') \ge n' - 6 = q' - 1$.  Consequently,
$B^{(\lambda')}_{q'+3} V_{\lambda'} \subseteq E_j^-\!\mid_{V_{\lambda'}}$
for $1 \le j \le q' - 2$, and in particular $B^{(\lambda')}_{q'+3} \subseteq E_1^-$
for $q' \ge 3$.
\end{theorem}

\begin{theorem}[Extended sign-kill at $\hatl = (4, 2)$, May-14 sixth pass]\label{thm:42-prefix}
For $\lambda' = (4, 2, 1^{q'}) \vdash n' = q' + 6$ with $q' \ge 4$,
every $T' \in \supp\bigl(B^{(\lambda')}_{q'+3} V_{\lambda'}\bigr)$
satisfies $p(T') \ge n' - 8 = q' - 2$.  Consequently,
$B^{(\lambda')}_{q'+3} V_{\lambda'} \subseteq E_j^-\!\mid_{V_{\lambda'}}$
for $1 \le j \le q' - 3$, and in particular $B^{(\lambda')}_{q'+3} \subseteq E_1^-$
for $q' \ge 4$.
\end{theorem}

\begin{remark}\label{rem:42-gap}
Theorem~\ref{thm:42-prefix} as stated in the sixth-pass paper has a mild
gap (its $E_{n'-1}^+$ decomposition silently omits the same-row $+q$-piece
at $c_1 = (1, 4)$ with inner shape $(2, 2, 1^{q'})$).  The eighth-pass
paper (Remark~3.7 there) and the present paper note that this gap is
closable by adding an extra branch handled by the $\hatl = (2, 2)$
extended sign-kill theorem (May-14 fourth pass).  The conclusion of
Theorem~\ref{thm:42-prefix} (containment direction) is correct as
stated; we apply it below.
\end{remark}

\begin{theorem}[Hook column-$1$, May-12]\label{thm:hook-E1minus}
For $\mu = (k, 1^m) \vdash n_0$ with $k \ge 2$ and $m \ge \max(k, 2)$,
$B^{(\mu)}_{\ell_0+1} V_\mu \subseteq E_1^-\!\mid_{V_\mu}$ (where
$\ell_0 = \ell(\mu) = m + 1$).
\end{theorem}

\begin{theorem}[Hook support lemma, May-14 fifth pass]\label{thm:hook-support}
Let $\mu = (k, 1^m) \vdash n_0 := k + m$ with $k \ge 2$ and $m \ge k$.
Set $\ell_0 := \ell(\mu) = m + 1$ and $j_0 := \ell_0 + 1 = m + 2$.
Then every $T \in \supp(B^{(\mu)}_{j_0} V_\mu)$ satisfies
$p(T) \ge m - k + 2$.
\end{theorem}

\begin{theorem}[Phase A endpoint, May-19]\label{thm:phaseA}
For any $\Hq$-module $V$ and any $1 \le j \le n-1$,
$(T_j+1)\cdots(T_1+1) V = E_j^+(V)$.  In particular,
$\Rp_n V_\lambda = E_{n-1}^+|_{V_\lambda}$.
\end{theorem}

\begin{lemma}[Outer-factor preservation, May-14 fifth pass]\label{lem:outer-preserve}
Let $W \subseteq V_\nu$ be a subspace such that every $T \in \supp(W)$
satisfies $p(T) \ge P$.  Then for every $j \ge P + 1$,
every $T'' \in \supp((T_j{+}1) W)$ also satisfies $p(T'') \ge P$.
\end{lemma}

\section{The four-piece $E^+_{n-1}$ decomposition for $(5,2,1^q)$}\label{sec:Eplus}

The $+q$-eigenspace $E_{n-1}^+|_{V_\lambda}$ admits a clean
four-piece decomposition reflecting the corner structure of $\lambda$.

\paragraph{Corner-pair pieces.}  For each pair $\{c_X, c_Y\}$ of distinct
removable corners, placement of $\{n-1, n\}$ at $\{c_X, c_Y\}$ in either
order yields a $T_{n-1}$-$2$-block in $V_\lambda$.  The pairs and axial
distances:

\begin{itemize}[topsep=0.2em,itemsep=0.1em]
\item $\{c_1, c_2\} = \{(1, 5), (2, 2)\}$.  Content $c(1, 5) - c(2, 2)
   = 4 - 0 = 4$.  Non-degenerate (\(|d| = 4 \ge 2\)).  Inner shape
   $\mu_C := (4, 1^{q+1})$.
\item $\{c_1, c_3\} = \{(1, 5), (q+2, 1)\}$.  Content $c(1, 5) -
   c(q+2, 1) = 4 - (1 - (q+1)) = q + 4$.  Non-degenerate (\(|d| \ge 9\)
   at \(q \ge 5\)).  Inner shape $\mu_A := (4, 2, 1^{q-1})$.
\item $\{c_2, c_3\} = \{(2, 2), (q+2, 1)\}$.  Content $c(2, 2) -
   c(q+2, 1) = 0 - (1 - (q+1)) = q$.  Non-degenerate (\(|d| \ge 5\) at
   \(q \ge 5\)).  Inner shape $\mu_B := (5, 1^q)$.
\end{itemize}

(The corner-pair labels $A, B, C$ match those of the May-13 night
$(5, 2)$ $j$-formula paper; in the eighth-pass $(4, 3)$ notation these
would be $\mu_{13}, \mu_{23}, \mu_{12}$, respectively.)

\paragraph{Same-row piece.}  A same-row $+q$-piece exists for $T_{n-1}$
when there is a removable corner $c$ and an SYT of $\lambda$ with $n-1$
at the cell immediately left of $c$ (same row) and $n$ at $c$.  The
candidates:

\begin{itemize}[topsep=0.2em,itemsep=0.1em]
\item At $c_1 = (1, 5)$: requires $T(1, 4) = n - 1$, $T(1, 5) = n$.
   Columns $4$ and $5$ contain only cells $(1, 4)$ and $(1, 5)$
   respectively (since $\lambda_2 = 2$), so no column-strict constraint
   below them.  The remaining $n - 2$ entries fill $\lambda \setminus
   \{(1, 4), (1, 5)\}$, which is the shape $\mu_h := (3, 2, 1^q)$, a
   valid partition.  \emph{Same-row piece at $c_1$ with inner shape
   $\mu_h$.}
\item At $c_2 = (2, 2)$: requires $T(2, 1) = n - 1$, $T(2, 2) = n$.
   But then column-$1$ strict increase forces $T(3, 1) > T(2, 1) = n - 1$,
   so $T(3, 1) = n$; conflict with $T(2, 2) = n$ (at $q \ge 1$).
   \emph{No same-row piece at $c_2$.}
\item At $c_3 = (q+2, 1)$: no cell to the left of column $1$.
   \emph{No same-row piece at $c_3$.}
\end{itemize}

\paragraph{No other $+q$-pieces.}  Same-column pairs at column $1$ give
$-1$-eigenvectors of $T_{n-1}$ (not in $E_{n-1}^+$).

\begin{lemma}[Four-piece decomposition of $E_{n-1}^+$]\label{lem:Eplus-decomp}
Define embeddings $\Phi_X : V_{\mu_X} \hookrightarrow V_\lambda$ for
$X \in \{A, B, C\}$ on the seminormal basis by
\[
   \Phi_X(v_{T'}) \;:=\; v_{T^X_+} + \tau_X\,v_{T^X_-},
\]
where $T^X_+, T^X_-$ are the two SYTs of $\lambda$ obtained from $T'$ by
placing $\{n-1, n\}$ at the corresponding corner pair in the two
possible orders, and $\tau_X \in \mathbb{Q}(q)^\times$ is the Hoefsmit
2-block coefficient at the corresponding axial distance.  Define
$\Phi_h : V_{\mu_h} \hookrightarrow V_\lambda$ on the seminormal basis
by $\Phi_h(v_{T''}) := v_{T'''}$, where $T'''$ is the unique SYT of
$\lambda$ obtained from $T''$ by placing $n - 1$ at $(1, 4)$ and $n$ at
$(1, 5)$.  Then:

\begin{enumerate}[topsep=0.3em,itemsep=0.1em,leftmargin=2em,
                  label=(\roman*)]
\item Each $\Phi_X$ ($X \in \{A, B, C, h\}$) is injective and
   $T_i$-equivariant for $1 \le i \le n - 3$.
\item The four images in $V_\lambda$ have pairwise disjoint
   seminormal supports and satisfy
   \[
       \Rp_n V_\lambda \;=\; E_{n-1}^+|_{V_\lambda}
       \;=\; \Phi_A(V_{\mu_A}) \;\oplus\; \Phi_B(V_{\mu_B})
       \;\oplus\; \Phi_C(V_{\mu_C}) \;\oplus\; \Phi_h(V_{\mu_h}).
   \]
\end{enumerate}
\end{lemma}

\begin{proof}
(i) Injectivity and $T_i$-equivariance for $i \le n - 3$ are standard
(May-20 Lemma 3.2; May-14 seventh pass Lemma~3.2): for $i \le n - 3$,
the swap $s_i$ moves only entries $i, i+1 \le n - 2$, both inside the
``inner'' cells of $\mu_X$ (or $\mu_h$); Hoefsmit coefficients depend
only on inner cells, so $T_i$ commutes with the corresponding $\Phi$.

(ii) Pairwise disjointness: each $\Phi$'s image is supported on SYTs of
$\lambda$ where $\{n - 1, n\}$ occupies a specific \emph{unordered} pair
of cells:
\begin{itemize}[topsep=0.2em,itemsep=0.1em]
\item $\Phi_A$: $\{(1, 5), (q+2, 1)\}$.
\item $\Phi_B$: $\{(2, 2), (q+2, 1)\}$.
\item $\Phi_C$: $\{(1, 5), (2, 2)\}$.
\item $\Phi_h$: $\{(1, 4), (1, 5)\}$ (specifically $n - 1$ at $(1, 4)$,
   $n$ at $(1, 5)$).
\end{itemize}
The four unordered cell pairs are pairwise distinct, so the supports
are pairwise disjoint.

Dimension match: $\dim \Phi_X(V_{\mu_X}) = f^{\mu_X}$ (injective), and
the $+q$-eigenvectors of $T_{n-1}$ in $V_\lambda$ partition into the
above four classes by where $\{n - 1, n\}$ sit.  By
Theorem~\ref{thm:phaseA}, $\Rp_n V_\lambda = E_{n-1}^+|_{V_\lambda}$, so
the direct-sum decomposition follows.

[Computationally verified at $q \in \{3, 4, 5\}$: $\dim \Rp_n V_\lambda$
equals $f^{\mu_A} + f^{\mu_B} + f^{\mu_C} + f^{\mu_h}$.  See
\texttt{\textasciitilde/projects/scratch/2026-05-14-extended-sign-kill-52/verify\_decomp.py}.]
\qed
\end{proof}

\begin{remark}\label{rem:two-hooks}
For $\hatl = (5, 2)$, two of the four inner shapes are hooks
($\mu_B = (5, 1^q)$ and $\mu_C = (4, 1^{q+1})$).  Compare with $\hatl =
(4, 3)$ (eighth pass), where only one of the four was a hook
($\mu_h = (4, 1^{q+1})$).  The shape-dependent feature is: in $(5, 2)$,
$\lambda_2 = 2$ is small enough that the same-row piece at $c_1$
strips two cells from row~1 leaving $(3, 2, 1^q)$ (non-hook).  In
$(4, 3)$, $\lambda_2 = 3$ is large enough that the analogous same-row
strip at $c_1$ would force $T(2, 3) > n - 1$, killing that branch ---
but the same-row at $c_2 = (2, 3)$ strips two cells from row~2 leaving
$(4, 1^{q+1})$ (hook).  The total piece count is the same (four), but
the hook count flips.
\end{remark}

\section{The two-branch reduction}\label{sec:reduction}

\begin{theorem}[Two-branch reduction]\label{thm:reduction}
For $\lambda = (5, 2, 1^q)$ with $q \ge 3$,
\begin{equation}\label{eq:reduction}
   B \;=\; (T_{n-5}{+}1)(T_{n-4}{+}1)(T_{n-3}{+}1)(T_{n-2}{+}1)\,
   \Bigl[\,\Phi_A\bigl(B_{\ell(\mu_A)+1}^{(\mu_A)} V_{\mu_A}\bigr)
        + \Phi_B\bigl(B_{\ell(\mu_B)+1}^{(\mu_B)} V_{\mu_B}\bigr)
   \,\Bigr].
\end{equation}
\end{theorem}

\begin{proof}
We follow the May-15/May-20 pull-through, applied separately to each
of the four pieces of $\Rp_n V_\lambda$ from
Lemma~\ref{lem:Eplus-decomp}, then show two pieces vanish.

\medskip\textbf{Step 1 (pulling four $\Rp$-blocks through $\Phi$).}
Let $\Phi \in \{\Phi_A, \Phi_B, \Phi_C, \Phi_h\}$ and $\mu$ the
corresponding inner shape.

Decompose $\Rp_{n-1} = (T_{n-2}+1)\,U_1$ where $U_1 = (T_{n-3}+1)\cdots
(T_1+1) = \Rp_{n-2}^{(\mu)} \in \Hq(S_{n-2})$.  Since $T_1, \ldots,
T_{n-3}$ commute with $T_{n-1}$ (index gaps $\ge 2$), $U_1$ preserves
$E_{n-1}^+|_{V_\lambda}$.  By Lemma~\ref{lem:Eplus-decomp}(i), $U_1
\Phi(v) = \Phi(\Rp_{n-2}^{(\mu)} v)$ for $v \in V_\mu$.  Hence
\[
   \Rp_{n-1}\Phi(V_\mu)
   \;=\; (T_{n-2}+1)\,\Phi\bigl(\Rp_{n-2}^{(\mu)} V_\mu\bigr).
\]

Iterating the same decomposition for $\Rp_{n-2}, \Rp_{n-3}, \Rp_{n-4}$,
using that each pulled-out $T_j$ factor commutes with the next inner
$\Rp$ (index gaps $\ge 2$):
\begin{equation}\label{eq:branchpull}
   \Rp_{n-4}\Rp_{n-3}\Rp_{n-2}\Rp_{n-1}\Phi(V_\mu)
   \;=\; (T_{n-5}+1)(T_{n-4}+1)(T_{n-3}+1)(T_{n-2}+1)\,
   \Phi\bigl(Q_\mu V_\mu\bigr),
\end{equation}
where
\[
   Q_\mu \;:=\; \Rp_{n-5}^{(\mu)}\, \Rp_{n-4}^{(\mu)}\, \Rp_{n-3}^{(\mu)}\, \Rp_{n-2}^{(\mu)}.
\]

\medskip\textbf{Step 2 (combining and applying Phase A).}
By Lemma~\ref{lem:Eplus-decomp}(ii) and Theorem~\ref{thm:phaseA},
\begin{align*}
   B \;=\; \Om V_\lambda
   &\;=\; \Rp_{n-4}\Rp_{n-3}\Rp_{n-2}\Rp_{n-1} \Rp_n V_\lambda \\
   &\;=\; \Rp_{n-4}\Rp_{n-3}\Rp_{n-2}\Rp_{n-1}
      \!\!\!\bigoplus_{X \in \{A,B,C,h\}}\!\!\!\Phi_X(V_{\mu_X}) \\
   &\;=\; (T_{n-5}{+}1)(T_{n-4}{+}1)(T_{n-3}{+}1)(T_{n-2}{+}1)\,
        \sum_{X}\Phi_X\bigl(I_X\bigr),
\end{align*}
where $I_X := Q_{\mu_X} V_{\mu_X}$.

\medskip\textbf{Step 3 (the $\mu_C$ piece vanishes).}
$\mu_C = (4, 1^{q+1})$ has $\ell(\mu_C) = q + 2 = n - 5$, hence
$\ell(\mu_C) + 1 = n - 4$.  The inner quadruple product is
\[
   I_C \;=\; \Rp_{n-5}^{(\mu_C)}\,\bigl(\Rp_{n-4}^{(\mu_C)}\,
   \Rp_{n-3}^{(\mu_C)}\,\Rp_{n-2}^{(\mu_C)} V_{\mu_C}\bigr)
   \;=\; \Rp_{n-5}^{(\mu_C)}\,B^{(\mu_C)}_{n-4} V_{\mu_C}.
\]
By Theorem~\ref{thm:hook-E1minus} applied to $\mu_C = (4, 1^{q+1})$
with $k = 4$, $m = q + 1$ (hypothesis $m \ge \max(k, 2) = 4$, i.e.,
$q \ge 3$ $\checkmark$), $B^{(\mu_C)}_{n-4} V_{\mu_C} \subseteq
E_1^-|_{V_{\mu_C}}$.  The rightmost factor of $\Rp_{n-5}^{(\mu_C)}$ is
$(T_1+1)$, which annihilates $E_1^-|_{V_{\mu_C}}$.  Hence $I_C = 0$.

\medskip\textbf{Step 4 (the $\mu_h$ piece vanishes).}
$\mu_h = (3, 2, 1^q)$ has $\ell(\mu_h) = q + 2 = n - 5$, hence
$\ell(\mu_h) + 1 = n - 4$.  Same argument as Step~3:
\[
   I_h \;=\; \Rp_{n-5}^{(\mu_h)}\,B^{(\mu_h)}_{n-4} V_{\mu_h}.
\]
By Theorem~\ref{thm:32-prefix} applied to $\mu_h = (3, 2, 1^q)$ with
$q' = q \ge 3$ (joint hypothesis of the theorem and the present proof;
we will tighten to $q \ge 5$ below), $B^{(\mu_h)}_{n-4} V_{\mu_h}
\subseteq E_1^-|_{V_{\mu_h}}$ (the $j = 1$ instance of
Theorem~\ref{thm:32-prefix}'s chain).  The rightmost $(T_1+1)$ of
$\Rp_{n-5}^{(\mu_h)}$ kills this, so $I_h = 0$.

\medskip\textbf{Step 5 (the $\mu_A$ and $\mu_B$ pieces are standard
$B^{(\mu_X)}$).}  Both $\mu_A = (4, 2, 1^{q-1})$ and $\mu_B = (5, 1^q)$
have $\ell(\mu_X) = q + 1 = n - 6$, hence $\ell(\mu_X) + 1 = n - 5$.
The inner quadruple product runs from $\Rp^{(\mu_X)}_{n-5}$ to
$\Rp^{(\mu_X)}_{n-2}$, exactly the four blocks of $\Om^{(\mu_X)}$.
Therefore
\[
   I_A \;=\; B^{(\mu_A)}_{\ell(\mu_A)+1} V_{\mu_A},
   \qquad
   I_B \;=\; B^{(\mu_B)}_{\ell(\mu_B)+1} V_{\mu_B}.
\]

Combining Steps 2--5 gives~\eqref{eq:reduction}.\qed
\end{proof}

\begin{remark}\label{rem:dim-add}
The dimensions match:
\[
   \dim B \;=\; f^{(4, 1)} \;=\; 4,
   \qquad
   \dim B^{(\mu_A)}_{\ell+1} \;=\; f^{(3, 1)} \;=\; 3,
   \qquad
   \dim B^{(\mu_B)}_{\ell+1} \;=\; f^{(4)} \;=\; 1,
\]
and $4 = 3 + 1$ confirms that both contributing Phi-pieces inject and
the four outer $(T_j+1)$ factors are jointly injective on the sum
$\Phi_A(I_A) + \Phi_B(I_B)$ (computationally verified at $q = 5$).
\end{remark}

\section{Main theorem}\label{sec:main}

\begin{theorem}[Main]\label{thm:main}
For $\lambda = (5, 2, 1^q) \vdash n = q + 7$ with $q \ge 5$,
\[
   B \;\subseteq\; \bigcap_{j=1}^{q-4} E_j^-\!\mid_{V_\lambda}.
\]
\end{theorem}

\begin{proof}
By Remark~\ref{rem:prefix-implies-Ej-}, it suffices to show that every
$T \in \supp(B)$ satisfies $p(T) \ge q - 3 = n - 10$.  We use the
two-branch reduction~\eqref{eq:reduction} together with prefix bounds
on the two inner $B^{(\mu_X)}_{\ell+1}$ pieces and
Lemma~\ref{lem:outer-preserve} on the four outer factors.

\medskip\textbf{Step 1 (prefix bound on the $\Phi_A$-image).}
$\mu_A = (4, 2, 1^{q-1})$ has size $n - 2$, $\ell(\mu_A) = q + 1$, and
trailing-$1$-count $q' = q - 1$.  Since $q \ge 5$ we have $q' \ge 4$,
so Theorem~\ref{thm:42-prefix} applies, giving
\[
   p(T') \;\ge\; q' - 2 \;=\; q - 3
   \qquad \text{for every } T' \in \supp\bigl(B^{(\mu_A)}_{\ell(\mu_A)+1} V_{\mu_A}\bigr).
\]

The Phi-embedding $\Phi_A$ adds $\{n - 1, n\}$ at $\{c_1, c_3\} =
\{(1, 5), (q+2, 1)\}$.  The added cells: $(1, 5)$ sits in column $5$
(not column $1$); $(q+2, 1)$ sits at row $q + 2 = n - 5 > q + 1 =
\ell(\mu_A)$, i.e., \emph{below} all column-$1$ cells of $\mu_A$.  So
column-$1$ rows $1, \ldots, q + 1$ of any extension SYT $T^A_\pm$
match the column-$1$ entries of $T'$, and $p(T^A_\pm) \ge p(T') \ge
q - 3$.

\medskip\textbf{Step 2 (prefix bound on the $\Phi_B$-image).}
$\mu_B = (5, 1^q)$, a hook with $k = 5$, $m = q$.  Since $q \ge 5$ we
have $m \ge k = 5$, so Theorem~\ref{thm:hook-support} applies and
gives
\[
   p(T') \;\ge\; m - k + 2 \;=\; q - 3
   \qquad \text{for every } T' \in \supp\bigl(B^{(\mu_B)}_{\ell(\mu_B)+1} V_{\mu_B}\bigr).
\]

The Phi-embedding $\Phi_B$ adds $\{n - 1, n\}$ at $\{c_2, c_3\} =
\{(2, 2), (q+2, 1)\}$.  Cell $(2, 2)$ is column $2$ (not column $1$);
cell $(q+2, 1)$ is row $q + 2 = n - 5 > q + 1 = \ell(\mu_B)$, below
$\mu_B$'s column-$1$ cells.  So by the same argument as Step~1,
$p(T^B_\pm) \ge p(T') \ge q - 3$.

\medskip\textbf{Step 3 (outer factors preserve the prefix).}
The four outer factors in~\eqref{eq:reduction} are $(T_j{+}1)$ for
$j \in \{n - 5, n - 4, n - 3, n - 2\}$.  For prefix bound $P = q - 3
= n - 10$, the hypothesis $j \ge P + 1 = q - 2 = n - 9$ of
Lemma~\ref{lem:outer-preserve} requires $j \ge n - 9$.  The smallest
index $j = n - 5$ satisfies $n - 5 \ge n - 9$ iff $9 \ge 5$
$\checkmark$, true unconditionally.  Applying
Lemma~\ref{lem:outer-preserve} four times in sequence, the prefix
bound $\ge q - 3$ is preserved.

\medskip\textbf{Conclusion.}  By Steps 1--3 and~\eqref{eq:reduction},
every $T \in \supp(B)$ has $p(T) \ge q - 3$.  By
Remark~\ref{rem:prefix-implies-Ej-}, $B \subseteq E_j^-$ for $1 \le
j \le q - 4 = \tau(\lambda)$.\qed
\end{proof}

\begin{remark}[Hypotheses]\label{rem:hyp}
The hypothesis $q \ge 5$ is the joint requirement of:
\begin{itemize}[topsep=0.2em,itemsep=0.1em]
\item Theorem~\ref{thm:42-prefix} at $q' = q - 1 \ge 4$, i.e., $q \ge 5$
   (for $\mu_A$).
\item Theorem~\ref{thm:hook-support} at $\mu_B = (5, 1^q)$ with
   $k = 5$, $m = q \ge 5$ (for $\mu_B$).
\item Theorem~\ref{thm:32-prefix} at $\mu_h = (3, 2, 1^{q})$ with
   $q' = q \ge 3$ (for $\mu_h$ vanishing in Step~4 of
   Theorem~\ref{thm:reduction}).
\item Theorem~\ref{thm:hook-E1minus} at $\mu_C = (4, 1^{q+1})$ with
   $k = 4$, $m = q + 1 \ge 4$, i.e., $q \ge 3$ (for $\mu_C$ vanishing
   in Step~3).
\end{itemize}
The threshold for the two contributing-branch theorems ($q \ge 5$) is
tight: at $q = 4$ the prediction $\tau = q - 4 = 0$ is vacuous, and at
$q \le 3$ the prediction is negative.  Thus the hypothesis $q \ge 5$
matches the first $q$ at which $\tau \ge 1$.
\end{remark}

\section{Sharpness (computational)}\label{sec:sharp}

The predicted threshold $\tau(\hatl) = q + 3 - \hatl_1 - \hatl_2 =
q - 4$ asserts both
$B \subseteq \bigcap_{j=1}^{q-4} E_j^-$ (Theorem~\ref{thm:main}) and
$B \not\subseteq E_{q-3}^-$.  We verify the latter computationally
for the first non-vacuous case $q = 5$ (where $q - 3 = 2$).

\begin{theorem}[Sharpness at $q = 5$, computational]\label{thm:sharp-q5}
For $\lambda = (5, 2, 1^5)$ at $n = 12$, there exists $T \in \supp(B)$
with $p(T) = 2 = q - 3$ exactly.  Consequently $B \not\subseteq
E_{q-3}^-|_{V_\lambda} = E_2^-|_{V_\lambda}$.
\end{theorem}

\begin{proof}[Computational evidence]
Applying $\Om$ to the first $40$ Hoefsmit basis vectors and recording
the support (Section~\ref{sec:verify}) gives a probe set in $\supp(B)$
of size $524$, with prefix histogram
$\{2: 168,\ 3: 168,\ 4: 112,\ 5: 56,\ 6: 16,\ 7: 4\}$.  An explicit
witness at $p(T) = 2$:
\[
   T(1, \cdot) = (1, 3, 4, 5, 6),
   \quad T(2, \cdot) = (2, 7),
   \quad T(3..7, 1) = (8, 9, 10, 11, 12).
\]
Column-$1$ entries are $(1, 2, 8, 9, 10, 11, 12)$, so $p(T) = 2$
exactly.  Hence $T \in \supp(B)$ has $p(T) < q - 2$, and by the
contrapositive of Remark~\ref{rem:prefix-implies-Ej-} (applied at
$j = q - 3 = 2$), $B \not\subseteq E_2^-$.
\qed
\end{proof}

\begin{remark}\label{rem:structural-sharp}
A structural-sharpness argument analogous to the May-14 sixth-pass
$(4, 2)$ Lemma~5.1 should locate entry $n - 8$ at $(1, 3)$ in a forced
position and produce an off-diagonal $T_{q - 3}$-swap to a tableau
violating the prefix bound, contradicting $B \subseteq E_{q-3}^-$.
We omit the detailed write-up of this for the present paper; the
mechanism is identical to the $(4, 2)$ proof.
\end{remark}

\section{Computational verification}\label{sec:verify}

We verify Theorem~\ref{thm:main} and the structural decomposition at
$q = 5$ (so $n = 12$, $\dim V_\lambda = 1728$).

\paragraph{Decomposition.}  At $q = 5$ ($n = 12$, $\dim V_\lambda = 1728$),
\begin{align*}
   f^{\mu_A} &= f^{(4, 2, 1^4)} = 350, &
   f^{\mu_B} &= f^{(5, 1^5)} = 126, \\
   f^{\mu_C} &= f^{(4, 1^6)} = 84, &
   f^{\mu_h} &= f^{(3, 2, 1^5)} = 160.
\end{align*}
Sum $= 720 = \dim \Rp_n V_\lambda$, matching
Lemma~\ref{lem:Eplus-decomp}.  (Also verified at $q = 3$:
$90 + 35 + 35 + 64 = 224$; $q = 4$: $189 + 70 + 56 + 105 = 420$.)

\paragraph{$E_1^-$ containment at $q = 5$.}  Apply $\Om$ to five
independent random vectors in $V_\lambda$ (mod $P = 100\,003$ at
$q = 11$) and check $(T_1 + 1)\,\Om v = 0$ on each.  All five tests
pass; each $\Om v$ has $524$ nonzero entries, and each $(T_1 + 1)\,\Om
v$ is the zero vector.  Probing $\dim B$ via $\Om$ applied to the
first $40$ basis vectors gives rank $4$, matching $f^{(4, 1)} = 4$.

\paragraph{Prefix bound and sharpness witness.}  The same probe (see
Theorem~\ref{thm:sharp-q5}) gives prefix histogram
$\{2, 3, 4, 5, 6, 7\}$ with min $= 2 = q - 3$, matching the predicted
threshold and confirming sharpness at $q = 5$.

Scripts:
\texttt{\textasciitilde/projects/scratch/2026-05-14-extended-sign-kill-52/verify\_decomp.py,
verify\_q5\_sparse.py}.

\section{Updated table of proved $\hatl$ rows}\label{sec:table}

\begin{table}[h]
\centering
\renewcommand{\arraystretch}{1.15}
\caption{Status of the extended sign-kill conjecture by $\hatl$.  The
$\tau$ column lists the predicted threshold $\tau(\hatl) = q + 3 -
\hatl_1 - \hatl_2$.}
\begin{tabular}{l|c|l|l}
\toprule
$\hatl$ & predicted $\tau$ & status & paper \\
\midrule
$(2, 2)$ & $q - 1$ & proved $q \ge 2$           & May-14 fourth pass \\
$(3, 2)$ & $q - 2$ & proved $q \ge 3$           & May-14 fifth pass \\
$(4, 2)$ & $q - 3$ & proved $q \ge 4$ (gap closable) & May-14 sixth pass \\
$(3, 3)$ & $q - 3$ & proved $q \ge 4$           & May-14 seventh pass \\
$(4, 3)$ & $q - 4$ & proved $q \ge 5$            & May-14 eighth pass \\
$(4, 4)$ & $q - 5$ & proved $q \ge 6$            & May-14 ninth pass \\
$(5, 2)$ & $q - 4$ & \textbf{proved $q \ge 5$ (this paper)} & May-14 tenth pass \\
\bottomrule
\end{tabular}
\end{table}

\section{Discussion}\label{sec:discussion}

\subsection{The $\hatl_1 + \hatl_2 = 7$ row is now complete}

The eighth-pass $(4, 3)$ paper proved the first member of the
$\hatl_1 + \hatl_2 = 7$ row.  This paper proves the second and last
member, $(5, 2)$.  The two share predicted threshold $\tau = q - 4$
and the same hypothesis $q \ge 5$.  Combined with $(4, 4)$ at
$\tau = q - 5$ (ninth pass), the $|\hatl| \le 8$ stratum is closed
except for $(5, 3)$ at $|\hatl| = 8$.

\subsection{Hooks as contributing branches}

In $\hatl = (5, 2)$, the contributing branch $\mu_B = (5, 1^q)$ is a
hook.  This is the first proof in the May-14 series where a hook gives
a non-zero contribution to $B$ at the $\ell + 1$ index (rather than
vanishing as $\mu_C$ does).  The mechanism: the hook $\mu_B$ has
$\ell(\mu_B) = q + 1 = n - 6$, so the inner $\Om^{(\mu_B)}$ has four
blocks matching the four ``standard-position'' blocks of the
two-branch reduction, and the hook support lemma at $k = 5$, $m = q$
gives a clean prefix bound $\ge q - 3$.  The contribution
$B^{(\mu_B)}_{\ell+1} V_{\mu_B}$ is one-dimensional (hook dim), so
$\Phi_B$ contributes exactly $1$ to $\dim B$.

\subsection{Toward $\hatl = (5, 3)$}

$\lambda = (5, 3, 1^q)$ has corners $c_1 = (1, 5)$, $c_2 = (2, 3)$,
$c_3 = (q + 2, 1)$.  Three corner pairs and (possibly) one same-row
piece at $c_2$ give the same four-piece structure.  Inner shapes for
the corner pairs:
\begin{itemize}[topsep=0.2em,itemsep=0.1em]
\item $\mu_{12} = (4, 2, 1^q)$ --- $\hatl = (4, 2)$ shape.
\item $\mu_{13} = (4, 3, 1^{q-1})$ --- $\hatl = (4, 3)$ shape, recursive
   on eighth pass.
\item $\mu_{23} = (5, 2, 1^{q-1})$ --- $\hatl = (5, 2)$ shape, recursive
   on \emph{this} paper.
\end{itemize}
The same-row piece at $c_2 = (2, 3)$ would have inner shape
$\lambda \setminus \{(2, 2), (2, 3)\} = (5, 1^{q+1})$, a hook.  (The
candidate same-row at $c_1 = (1, 5)$ would have inner shape $(3, 3,
1^q)$, which requires the column-$4$ constraint $T(1, 4) < T(2, 4)$ to
not bite; cell $(2, 4)$ does not exist in $(5, 3, 1^q)$ since
$\lambda_2 = 3$, so there is no constraint and the candidate
\emph{does} contribute.  So $\hatl = (5, 3)$ may actually have
\emph{two} same-row pieces, giving five total pieces in
$E^+_{n-1}$.)  Investigation of $(5, 3)$ is left for a follow-up
session.

\subsection{Push status}

\textbf{71 unpushed commits} on \texttt{clio-vega/proofs} (after
today's tenth-pass commit, assuming the count from the ninth-pass paper
held at $70$ before this).  PAT issue from May~6 persists --- please
refresh.

\section{Status}\label{sec:status}

\paragraph{Established.}
\begin{itemize}[topsep=0.3em,itemsep=0.1em]
\item Lemma~\ref{lem:Eplus-decomp}: four-piece $E^+_{n-1}$
   decomposition for $\lambda = (5, 2, 1^q)$ at $q \ge 3$.
\item Theorem~\ref{thm:reduction}: two-branch reduction (two
   contributing, two vanishing) at $q \ge 3$.
\item Theorem~\ref{thm:main}: $B \subseteq E_j^-$ for $j \le q - 4$
   at $q \ge 5$.
\item Theorem~\ref{thm:sharp-q5}: sharpness at $j = q - 3$ for
   $q = 5$ (computational witness).
\end{itemize}

\paragraph{Open.}
\begin{itemize}[topsep=0.3em,itemsep=0.1em]
\item Structural (non-computational) sharpness across all $q \ge 5$
   (analogous to the sixth-pass Lemma~5.1; mechanism stated in
   Remark~\ref{rem:structural-sharp}).
\item Sharpness witnesses at $q \in \{6, 7, \ldots\}$ (consistent
   computationally up to memory limits; not run in this session).
\item The next family $\hatl = (5, 3)$ (potentially five-piece
   decomposition; section~\ref{sec:discussion}).
\end{itemize}

\end{document}
